\section{Introduction}
\enlargethispage{-7\baselineskip}
\label{md:introduccion}
Reading a constant as a scalar value represents the final stage of an
operation. Retaining the structure that produces this readout makes
further questions accessible: which relations it shares with other
readouts, which transformation preserves its orientation, and which part
of its history can be reconstructed from joint observations. The present
development situates these questions within the continuity between the
discrete HMT structure and its realizations in Dimensional Mechanics.

The common origin retains the operations of APP, the oriented regime of
TRIT, and the effective composition of TPK. Phase return is distinguished
from the evolution of the enriched state. The regional coordinates and
their dodecaphase closure are subsequently transmitted to the action,
incidence, and mass readout operators. Each transmission preserves the
objects and conditions it uses; an output can enter the next construction
without losing its provenance. The opening core presents this architecture
with its operations, prolongations, and proofs. The sectoral antecedents
that follow explicitly supply the dependencies used by the compositions
developed here.

The central question can be formulated through three complementary terms:
\emph{production}, \emph{compatibility}, and \emph{reconstruction}.
Production fixes the object and its genealogy. Compatibility determines
the relations that its different readouts must satisfy. Reconstruction
establishes which data of the object can be recovered from them. This
inverse procedure is a property of already constructed readouts and
retains their bases, orientations, and domains; it does not replace their
generative order.

The first line of development links action to the angular pair. The
dimensionless coefficient of the action section participates in the
counter-angle; together with the angle, the latter determines the channels
that enter the constitutive response and the Barbero--Immirzi functional.
The same section enters the return quotient of the electronic readout.
Substituting the structures into their equations therefore reveals shared
dependencies that a list of separate constants leaves implicit. The
distinction between a coefficient and a dimensionful quantity is
essential: reconstructing a shape ratio and reconstructing a dimensionful
action require different reference data.

The second line of development examines the information retained in the
response. Normalized speed and impedance retain different aspects of the
distribution between channels. Their composition proves both a global
inversion within the stated collection and a differential relation between
the two responses. A spectral potential organizes this relation, and its
Hessian controls the sensitivity of inversion. Two properties are thus
distinguished: exact data uniquely determine the object, whereas finite
precision permits reconstruction with a bounded error. The oriented
collection whose extremal value is Catalan's constant enters through its
spectral operation, with its arguments and depth retained.

Transported composition of the responses provides their additive
coordinate; the angular gauge and Higgs signatures reconstruct the same
channel pair. The moments of the cyclic defect complete this reading
with convergence bounds and a Mellin representation. Subsequently, the
electronic section is substituted into Fermi and the Higgs scale: the
proof retains the factors that remain, those that cancel, and the
transformation of the inverse mass operator without assuming commutation.

The third line of development addresses the order of operations. The
quarter-turn plane of the dodecaphase cycle and the conjugate forms of the
ellipse yield an explicit commutator. The balance of phase increments
closes while the action on the state retains a reciprocal dilation. Two
programs with the same multiplicities of operations can produce different
responses because of their order. This difference admits defined
detectors and reconstruction maps: the trace records certain invariants,
whereas responses along oriented axes distinguish directions to which the
trace assigns the same value.

The physical interpretation is presented together with the conditions of
each realization. In the circular gravitational section, the reduced
Compton wavelength and the reduced gravitational radius undergo
reciprocal dilations and retain a common area. In the thermal realization
of the cycle, the same structural coordinate enters as the exponent of an
equilibrium weight. In elliptic transport, preservation of symplectic area
and preservation of a common quadratic energy are examined through their
respective operators. Any prospective experimental realization of the relative
program requires a joint specification of the system, reference,
available operations, and controller. These conditions are part of the
meaning of the result and remain adjacent to its proofs.

Gravitational normalization is developed before its scalar consequences
are used. The Hadamard return, the inscription of incidences, and the
local norm of $\alpha$ produce a length section. Its polar transport
and the same route character produce two reciprocal readout operators.
Energy in turn preserves the action length. Combining
$R\Lambda=L_\alpha^2I$ with $H\Lambda=\hbar cI$ fixes the gravitational
constant as the common coefficient of all positive modes of the
realization. The circular clock is then expressed in terms of length
and speed. This precedence explains the normalization without using
the experimental value of $G$ or a tabulated Planck length.

The electronic and thermal compositions specify how memory enters these
relations. In the former, the dodecaphase record is reconstructed from
its sum, differences, and orientation; the spectrum distinguishes
orders having the same multiplicities. In the latter, the geometric
spectrum already constructed by memory transport supplies the ratio
$1/10$. Chronological capacities and occupation sectors determine a
marked trace that is composed with a Gibbs state. Conditional energy,
marked information, and the resulting temperature are specified
separately; this choice of observations makes it possible to prove the
thermometric transducer and follow its relation to the action scale and
area. The corresponding sections retain the preparation data, gradings,
and variation hypotheses used.

The exposition first develops the antecedents, then preserves the
explanatory narrative of the investigation and presents the complete
compositions. The chronological study of vacancies, the regional example
of visible recurrence, and memory re-entry make it possible to follow the
origin of the compatibility questions. The energetic and Gaussian
developments are retained in the appendices as auxiliary realizations,
with their full hypotheses and proofs. The conclusions bring together
the structural characterizations obtained. The digital supplement
preserves the verbatim working sources, programs, and results; the
manuscript contains the proofs used and distinguishes their mathematical
interpretation from the accompanying finite checks.

This organization gives a precise formulation of the work's transversal
contribution. A constant retains its numerical value, while its structural
definition also specifies the operation producing that value, the domain
on which it acts, and the transformations it preserves. Two quantities
sharing a dependency can therefore constrain one another: knowing a
response and its orientation reconstructs a distribution that an isolated
scalar readout does not distinguish. The inverse problem is posed on the
already generated image and retains its reference data. Generation and
reconstruction are thereby connected without reversing the precedence of
the construction.

The same distinction organizes the relation between memory and
information. The present state, the accumulated balance, and the operator
that will respond to a continuation are different objects. The return of
a phase can coexist with a persistent transformation of the state; two
sequences with identical counts can differ through their order of
composition. The proofs develop this distinction through operators,
invariants, and reconstructive observations. Their physical relevance
lies in specifying what each description preserves and which operation
provides access to the retained information.

The passage to dimensionful quantities preserves the corresponding bases
and normalizations. Reconstruction of an action coefficient, preservation
of an area, and characterization of a temperature are related results of
different types. The exposition brings them together under their
conditions of validity and shows how they compose. The conclusions make
explicit both the characterization obtained for each object and the
relation it establishes with the other structures.

The Feigenbaum constant acquires a precise structural definition: it is the asymptotic separation invariant of the parameters selected by the critical returns of the uniform dodecaphasic reader. The construction preserves its APP--TRIT--TPK origin, the joint continuum and the first-root order. Identifying the selected centres with the analytic centres makes the scaling theorem applicable to that same sequence; the proofs in \ref{hmtfg:escalado:15}--\ref{hmtfg:escalado:17} and \eqref{hmtfg:articulacion:limite} establish the parametric result. The Jacobi representation recovers the profile from its twelve-node measure, while constitutive inversion recovers the oriented channels and the moment whose endpoint is Catalan's constant. These operations exhibit relations between readers of the same antecedent, preserving the domain and information of each.

The composition also extends to a common action of geometry, gauge fields and matter. Gravitation and the strong, weak and electromagnetic interactions enter through realizations of the same enriched state, with the couplings inherited from the preceding construction. Joint variation determines the sources and torsional response; the Legendre transformation produces a common algebra of first-class constraints. Their propagation and the closure of the canonical connection are proved on the regular constraint surface, preserving boundary conditions and global holonomy. Variational and canonical unification thus provides an explicit dynamical composition for the compatibility and reconstruction relations organizing the manuscript.
